\documentclass[a4paper,11pt]{article}

\usepackage[italian]{babel}
\usepackage[T1]{fontenc}
\usepackage[utf8]{inputenc}
\usepackage{lmodern}
\usepackage{geometry}
\usepackage{xcolor}
\usepackage{hyperref}
\usepackage{tabularx}
\usepackage{array}
\usepackage{booktabs}
\usepackage{enumitem}
\usepackage{titlesec}
\usepackage{tcolorbox}
\usepackage{fancyhdr}
\usepackage{multicol}
\usepackage{amsmath}

\geometry{margin=2.2cm}
\setlength{\parskip}{0.6em}
\setlength{\parindent}{0pt}
\setlength{\headheight}{14pt}

% --- Colori ---
\definecolor{MainColor}{HTML}{1F4E79}
\definecolor{AccentColor}{HTML}{2E86AB}
\definecolor{SoftBlue}{HTML}{EAF4FB}
\definecolor{SoftGray}{HTML}{F5F7FA}
\definecolor{SoftGreen}{HTML}{EAF8F0}
\definecolor{DarkText}{HTML}{1E1E1E}

% --- Hyperref ---
\hypersetup{
    colorlinks=true,
    linkcolor=MainColor,
    urlcolor=AccentColor,
    pdftitle={Guida moderna a GitHub},
    pdfauthor={}
}

% --- Titoli ---
\titleformat{\section}
  {\Large\bfseries\color{MainColor}}
  {\thesection}{0.8em}{}

\titleformat{\subsection}
  {\large\bfseries\color{AccentColor}}
  {\thesubsection}{0.8em}{}

% --- Header/Footer ---
\pagestyle{fancy}
\fancyhf{}
\fancyhead[L]{\textcolor{MainColor}{Guida a GitHub}}
\fancyhead[R]{\textcolor{AccentColor}{Version control e collaborazione}}
\fancyfoot[C]{\thepage}

% --- Box ---
\newtcolorbox{infobox}{
    colback=SoftBlue,
    colframe=AccentColor,
    boxrule=0.8pt,
    arc=3mm,
    left=2mm,right=2mm,top=2mm,bottom=2mm
}

\newtcolorbox{tipbox}{
    colback=SoftGreen,
    colframe=green!50!black,
    boxrule=0.8pt,
    arc=3mm,
    left=2mm,right=2mm,top=2mm,bottom=2mm
}

\newtcolorbox{codebox}{
    colback=SoftGray,
    colframe=black!20,
    boxrule=0.5pt,
    arc=2mm,
    left=2mm,right=2mm,top=2mm,bottom=2mm
}

\begin{document}

% =========================
% Titolo
% =========================
\vspace*{1.5cm}

{\Huge\bfseries\color{MainColor} Guida a \texttt{GitHub}} \\[0.4cm]
{\Large\color{AccentColor} Dalle basi al workflow collaborativo completo}

\vspace{1cm}

\begin{infobox}
\textbf{GitHub} e la piattaforma dove pubblichi repository Git, collabori tramite pull request, tieni traccia del lavoro con issue e automatizzi con Actions.
\end{infobox}

\vspace{0.5cm}

\begin{codebox}
\textbf{Quick reference}

\renewcommand{\arraystretch}{1.2}
\begin{tabularx}{\textwidth}{>{\ttfamily}l X}
\toprule
Comando / Azione & Descrizione \\
\midrule
git clone URL & Clona repository in locale \\
git checkout -b feature/nome & Crea e apre un branch \\
git add . & Aggiunge modifiche allo staging \\
git commit -m "messaggio" & Crea un commit \\
git push -u origin feature/nome & Pubblica branch su GitHub \\
Pull Request & Richiesta di merge su GitHub \\
git pull --rebase origin main & Allinea il branch locale \\
git fetch --all --prune & Aggiorna i riferimenti remoti \\
Issue & Traccia bug, task e richieste \\
Actions & Esegue CI/CD automatizzata \\
\bottomrule
\end{tabularx}
\end{codebox}

\vspace{0.6cm}
\tableofcontents
\newpage

% =========================
\section{Cos'e GitHub}
\texttt{GitHub} e un servizio cloud basato su \texttt{Git}. Git e il motore di versionamento; GitHub aggiunge collaborazione e strumenti di team.

GitHub e utile quando:
\begin{itemize}[leftmargin=1.2cm]
    \item vuoi salvare e versionare codice in remoto;
    \item lavori in gruppo su branch separati;
    \item vuoi revisionare modifiche prima del merge;
    \item vuoi automatizzare test e deploy.
\end{itemize}

\begin{tipbox}
\textbf{Idea pratica:} usa Git in locale per sviluppare e GitHub per condividere, revisionare e integrare.
\end{tipbox}

% =========================
\section{I concetti fondamentali}
\begin{infobox}
\textbf{Repository} \\
Contenitore del progetto (codice, storia commit, file di configurazione).

\medskip
\textbf{Branch} \\
Linea di sviluppo separata, utile per feature o fix.

\medskip
\textbf{Commit} \\
Snapshot coerente delle modifiche.

\medskip
\textbf{Pull Request (PR)} \\
Richiesta di revisione e merge di un branch su un altro.

\medskip
\textbf{Issue} \\
Ticket per bug, task, idee o discussioni operative.
\end{infobox}

Struttura mentale:
\[
\text{Repository} \rightarrow \text{Branch} \rightarrow \text{Commit} \rightarrow \text{Pull Request}
\]

% =========================
\section{Setup iniziale}
\subsection{Configura Git}
\begin{codebox}
\texttt{git config --global user.name "Nome Cognome"} \\
\texttt{git config --global user.email "mail@example.com"}
\end{codebox}

\subsection{Genera una chiave SSH}
\begin{codebox}
\texttt{ssh-keygen -t ed25519 -C "mail@example.com"} \\
\texttt{cat \~/.ssh/id\_ed25519.pub}
\end{codebox}

Copia la chiave pubblica e aggiungila su GitHub in:
\begin{codebox}
\texttt{Settings > SSH and GPG keys}
\end{codebox}

\subsection{Test connessione}
\begin{codebox}
\texttt{ssh -T git@github.com}
\end{codebox}

% =========================
\section{Primo flusso di lavoro}
\subsection{Clona il repository}
\begin{codebox}
\texttt{git clone git@github.com:org/progetto.git} \\
\texttt{cd progetto}
\end{codebox}

\subsection{Crea un branch di lavoro}
\begin{codebox}
\texttt{git checkout -b feature/login}
\end{codebox}

\subsection{Lavora e salva}
\begin{codebox}
\texttt{git add .} \\
\texttt{git commit -m "feat: aggiunge login base"}
\end{codebox}

\subsection{Pubblica su GitHub}
\begin{codebox}
\texttt{git push -u origin feature/login}
\end{codebox}

Poi, su GitHub, apri una \textbf{Pull Request} verso \texttt{main}.

% =========================
\section{Pull Request fatta bene}
Una PR efficace riduce tempi di review e problemi in merge.

\begin{center}
\renewcommand{\arraystretch}{1.3}
\begin{tabularx}{\textwidth}{>{\ttfamily}l X}
\toprule
Elemento & Buona pratica \\
\midrule
Titolo & Breve, orientato al risultato \\
Descrizione & Cosa cambia, perche, impatti \\
Scope & Piccolo e focalizzato \\
Checklist & Test, lint, documentazione \\
Reviewers & Assegna le persone giuste \\
\bottomrule
\end{tabularx}
\end{center}

\begin{tipbox}
\textbf{Consiglio utile:} prima di aprire la PR, aggiorna il branch con \texttt{main} per ridurre conflitti.
\end{tipbox}

% =========================
\section{Aggiornare il branch senza caos}
\subsection{Allinea il locale}
\begin{codebox}
\texttt{git checkout main} \\
\texttt{git pull origin main}
\end{codebox}

\subsection{Porta le novita sul branch feature}
\begin{codebox}
\texttt{git checkout feature/login} \\
\texttt{git rebase main}
\end{codebox}

Se ci sono conflitti:
\begin{enumerate}[leftmargin=1.2cm]
    \item risolvi i file;
    \item esegui \texttt{git add <file>};
    \item continua con \texttt{git rebase --continue}.
\end{enumerate}

Al termine, aggiorna il remoto:
\begin{codebox}
\texttt{git push --force-with-lease}
\end{codebox}

% =========================
\section{Issue e organizzazione}
Le \textbf{Issue} sono il punto di ingresso del lavoro.

\begin{center}
\renewcommand{\arraystretch}{1.3}
\begin{tabularx}{\textwidth}{>{\ttfamily}l X}
\toprule
Strumento & Uso pratico \\
\midrule
Issue template & Standardizza bug report e task \\
Label & Classifica priorita e tipo lavoro \\
Milestone & Raggruppa task per release \\
Project board & Stato operativo (todo/doing/done) \\
\bottomrule
\end{tabularx}
\end{center}

\begin{infobox}
\textbf{Regola semplice:} ogni PR dovrebbe riferirsi a una issue, cosi il contesto resta chiaro nel tempo.
\end{infobox}

% =========================
\section{README minimo consigliato}
Un buon repository ha almeno:
\begin{itemize}[leftmargin=1.2cm]
    \item descrizione breve del progetto;
    \item prerequisiti;
    \item istruzioni di installazione e avvio;
    \item esempi d'uso;
    \item sezione contributi.
\end{itemize}

\begin{codebox}
\texttt{\# Nome progetto} \\
\texttt{Descrizione breve} \\
\texttt{\#\# Setup} \\
\texttt{\#\# Uso} \\
\texttt{\#\# Contribuire}
\end{codebox}

% =========================
\section{GitHub Actions in 60 secondi}
\texttt{GitHub Actions} serve per CI/CD: test automatici, lint, build e deploy.

File tipico:
\begin{codebox}
\texttt{.github/workflows/ci.yml}
\end{codebox}

Schema base:
\begin{itemize}[leftmargin=1.2cm]
    \item trigger su \texttt{push} e \texttt{pull\_request};
    \item job con checkout codice;
    \item install dipendenze;
    \item esecuzione test/lint.
\end{itemize}

% =========================
\section{Cheat sheet veloce}
\begin{center}
\renewcommand{\arraystretch}{1.25}
\begin{tabularx}{\textwidth}{>{\ttfamily}l X}
\toprule
Comando / Azione & Significato \\
\midrule
git clone URL & Clona repository \\
git status & Mostra stato working tree \\
git checkout -b nome & Nuovo branch \\
git add . & Prepara modifiche per commit \\
git commit -m "msg" & Salva snapshot \\
git push -u origin nome & Pubblica branch \\
git pull --rebase origin main & Aggiorna branch locale \\
git fetch --all --prune & Pulisce riferimenti remoti vecchi \\
git log --oneline --graph & Cronologia compatta \\
Open Pull Request & Chiede review e merge \\
Merge Pull Request & Integra branch in target \\
Close Issue & Chiude task completato \\
\bottomrule
\end{tabularx}
\end{center}

% =========================
\section{Errori comuni dei principianti}
\begin{infobox}
\textbf{1. Commit troppo grandi} \\
Rendono difficile capire, revieware e fare rollback.

\medskip
\textbf{2. Lavorare direttamente su main} \\
Meglio usare branch dedicati per ogni cambiamento.

\medskip
\textbf{3. Messaggi commit vaghi} \\
Evita "fix" o "update": scrivi cosa e perche.

\medskip
\textbf{4. Aprire PR senza contesto} \\
Descrivi sempre obiettivo, test fatti e impatti.
\end{infobox}

% =========================
\section{Workflow consigliato in team}
\begin{enumerate}[leftmargin=1.2cm]
    \item Apri issue con obiettivo chiaro.
    \item Crea branch da \texttt{main}.
    \item Fai commit piccoli e descrittivi.
    \item Pubblica branch e apri PR.
    \item Ricevi review e applica feedback.
    \item Esegui merge e chiudi issue.
\end{enumerate}

% =========================
\section{Conclusione}
\texttt{GitHub} diventa potente quando usi un flusso ordinato: branch puliti, PR piccole, review serie e automazioni CI.

\begin{tipbox}
\textbf{Da ricordare:}
\begin{itemize}[leftmargin=1.2cm]
    \item separa sempre il lavoro in branch;
    \item apri PR leggibili e testate;
    \item usa issue per tracciare decisioni e task;
    \item automatizza il controllo qualita con Actions.
\end{itemize}
\end{tipbox}

\vfill
\begin{center}
\textcolor{AccentColor}{\large Buon lavoro con \texttt{GitHub}}
\end{center}

\end{document}
